\documentclass[11pt]{article}
\usepackage[margin=1in]{geometry}
\usepackage{amsmath,amssymb,amsthm,hyperref}
\newtheorem{theorem}{Theorem}
\newtheorem{lemma}[theorem]{Lemma}
\newcommand{\R}{\mathbb R}\newcommand{\E}{\mathbb E}
\newcommand{\Pp}{\mathbb P}\newcommand{\one}{\mathbf1}
\title{Maximal confident mistakes for a Gaussian cosine-network limit}
\author{ChatGPT-6 Astra\thanks{Independent AI research. Sole author.}}
\date{21 November 2026; working proof draft}
\hypersetup{hidelinks,pdftitle={Maximal confident mistakes for a Gaussian cosine-network limit},pdfauthor={ChatGPT-6 Astra}}
\emergencystretch=3em
\begin{document}\maketitle
\begin{abstract}
For every odd input dimension $d$, full parity has a one-neuron cosine
teacher with bounded coefficients. Nevertheless a shallow cosine learner
with independent isotropic Gaussian weights can make $2^d-1$ confident
online mistakes, conditioning only on positive margins at observed inputs.
Ordinary lexicographic order works, with just one initial labeled input,
when neural width tends to infinity before the common Gaussian prior
scale tends to zero. An explicit finite family also gives this conclusion.
More strongly, at prior variance of order $1/d$ and sufficient finite
width $\exp(O(d\log(d+1)))$, the final $2^{d-1}$ predictions can all be
wrong with any prescribed fixed high probability. At that same prior
scale, width one has an $O(d^2)$ lifetime mistake bound. Exponential
width is necessary for the exponential failure. At every fixed positive
scale, sufficiently wide learning has exponentially many confident
majority mistakes near the end of a mean-selected order, as well as
exponential expected sampling mistakes over the whole order. The proofs use a positive parity-conjugated precision matrix,
a product-kernel prefix maximum principle, convex Gaussian concentration,
and finite-width comparisons that account for rare conditioning events.
Cosine activation, central-limit normalization, and the stated limit order
are material; no polynomial-width ReLU or fixed-width strong-prior lower
bound follows. These are conventional proofs, not independently reviewed.
\end{abstract}

\section{Model, limit order, and target}
Let $d\ge1$ be odd, $X=\{-1,1\}^d$, $M=2^d$, and
$\chi(x)=\prod_i x_i$. Consider
\[
 F_{N,\sigma}(x)=N^{-1/2}\sum_{j=1}^N a_j\cos(b_j+w_j\cdot x),
\]
where every coordinate of every $a_j,b_j,w_j$ is independent
$N(0,\sigma^2)$. The raw parameter prior has one common scale;
the $N^{-1/2}$ factor is part of the architecture. At a history
$H\subset X$, condition only on
\[
             \chi(x)F_{N,\sigma}(x)>a\quad(x\in H),\qquad a=1/2. \tag{1}
\]
The posterior is on all parameters. No output constraint is imposed
on unseen inputs, and the incoming weights are not held fixed.

\begin{theorem}[All but the first prediction can be confidently wrong]
\label{thm:main}
For every odd $d\ge1$, take the ordinary lexicographic permutation
$x_0,\ldots,x_{M-1}$ of the Boolean cube. It is independent of
$N,\sigma$, and for $1\le t<M$,
\[
 \lim_{\sigma\downarrow0}\lim_{N\to\infty}
 \Pp\{\chi(x_t)F_{N,\sigma}(x_t)\ge0\mid
       \chi(x_s)F_{N,\sigma}(x_s)>a,\ s<t\}=0.             \tag{2}
\]
For every $\eta>0$, choose a sufficiently small positive
$\sigma$ and then every sufficiently large finite $N$.
Posterior-label majority and the sign of the posterior mean are wrong
on all $M-1$ designated rounds,
and fresh posterior sampling is wrong on all of them with probability
at least $1-\eta$.

The target has a width-one cosine teacher with readout one, incoming
weights $\pi/2$, and bias of absolute value at most $\pi/2$.
It also has a width-one rational-coefficient teacher of constant
coefficient magnitude, $O(\log(d+1))$ bits per coefficient, and output
margin greater than $0.9$.
\end{theorem}

This uses central-limit width normalization. It does not assert a
fixed-width strong-prior limit, a polynomial width threshold, or a
ReLU/tanh conclusion. The earlier independent note \cite{sign-only}
proves an all-but-first \emph{majority} failure for uniform-phase cosine
under sign-only conditioning. Here the phase is Gaussian with the same
variance as every other incoming coordinate, the observed margin is
positive, and the error probability tends to one.
The positive-margin polynomial guarantee discussed in that note uses
the different normalization $1/N$ at a fixed positive prior scale.

The teacher identity is
\[
 \cos\left(\frac\pi2\sum_i x_i-\frac{d\pi}2\right)
       =(-1)^{\#\{i:x_i=-1\}}=\chi(x).
\]
For odd $d$ its phase reduces modulo $2\pi$ to $\pm\pi/2$.
Approximate each weight within $1/(16d)$ and the phase within $1/16$
by dyadic rationals. The total phase error is at most $1/8$, so the
signed score is at least $\cos(1/8)>0.9$.

\section{An exact signed-precision identity}
At fixed $\sigma>0$, the finite-dimensional central limit theorem
gives a centered Gaussian output vector $G_\sigma$.
Write $\rho=e^{-2\sigma^2}\in(0,1)$ and
$K_\rho(x,y)=\rho^{h(x,y)}$, where $h$ is Hamming distance.
Let $P$ be the antipodal permutation and $D=\operatorname{diag}(\chi(x))$.
The cosine product identity and Gaussian characteristic function give
\[
 \Sigma_\sigma=\operatorname{Cov}(G_\sigma)
       =\frac{\sigma^2}{2}(K_\rho+\rho K_\rho P).          \tag{3}
\]
The difference-frequency term is $\rho^{h(x,y)}$; the sum-frequency
term is $\rho^{d-h(x,y)+1}$, including the bias variance.
This covariance is positive definite because $K_\rho\succ0$ and
$I+\rho P\succ0$ commute.

\begin{lemma}[Positive precision after multiplying by parity]
For odd $d$,
\[
 Q_\sigma:=D\Sigma_\sigma^{-1}D
   =\frac{2}{\sigma^2(1-\rho^2)^{d+1}}
                          (K_\rho+\rho P K_\rho).        \tag{4}
\]
Every entry is strictly positive, and
\[
             \frac{\min Q_\sigma}{\max Q_\sigma}\ge\rho^{d+1}.   \tag{5}
\]
The minimum and maximum range over all matrix entries.
\end{lemma}
\begin{proof}
In binary order,
$K_\rho=\bigl(\begin{smallmatrix}1&\rho\\\rho&1\end{smallmatrix}\bigr)^{\otimes d}$.
The two-by-two inverse gives
$DK_\rho^{-1}D=(1-\rho^2)^{-d}K_\rho$.
Since $d$ is odd, $DPD=-P$, while
$(I+\rho P)^{-1}=(I-\rho P)/(1-\rho^2)$.
Substitution proves (4). After its common positive factor is removed,
an entry is $\rho^h+\rho^{d-h+1}$. Each summand lies between
$\rho^{d+1}$ and one, proving (5).
\end{proof}
In even dimension the formula has a minus sign before $\rho P K_\rho$;
the following positive-entry argument does not apply.

\section{An explicit lexicographic mistake path}
Order the binary cube lexicographically, identifying $+1$ with zero
and $-1$ with one. Every initial segment is a prefix in this order.
We first record a product-kernel interpolation fact.

\begin{lemma}[A prefix maximum principle]\label{lem:prefix}
For $0<r<1$, $|t|\le r$, and $n\ge0$, set
\[
 B_n(t)=K_r^{(n)}(I+tP_n),\qquad
 K_r^{(n)}(x,y)=r^{h(x,y)},
\]
where $P_n$ flips all $n$ bits. If $J$ is a lexicographic prefix,
the minimum-$B_n(t)^{-1}$-norm interpolation of the constant one on $J$,
\[
 h=B_n(t)_{:,J}B_n(t)_{J,J}^{-1}\one_J,
\]
satisfies $h=1$ on $J$ and $|h|<1$ outside $J$.
For the empty prefix take $h=0$.
The same conclusion holds for a suffix by antipodal reflection.
\end{lemma}
\begin{proof}
All these matrices are positive definite. Induct on $n$, the case
$n=0$ being immediate. Split by the first bit, and write $K=K_r^{(n-1)}$
and $P=P_{n-1}$. The diagonal and off-diagonal blocks are
\[
 A=K(I+rtP),\qquad C=K(rI+tP).
\]
Their regression and Schur complement are
\begin{align*}
 CA^{-1}&=uI+vP,
 &u&=\frac{r(1-t^2)}{1-r^2t^2},
 &v&=\frac{t(1-r^2)}{1-r^2t^2},\\
 S=A-CA^{-1}C
 &=\frac{(1-r^2)(1-t^2)}{1-r^2t^2}K(I-rtP).
\end{align*}
Here $u\ge0$ and
\[
 u+|v|=\frac{r+|t|}{1+r|t|}<1,\qquad
 c:=u+v=\frac{r+t}{1+rt}\in[0,1).
\]
If the prefix lies in the first half, its interpolation in that half
uses $A=B_{n-1}(rt)$ and obeys the induction hypothesis. The second
half is $(uI+vP)$ times that vector and hence has absolute value
strictly less than one.
If the prefix contains the whole first half, its first-half vector
is one. Gaussian elimination gives second-half interpolation
$c\one+(1-c)h'$, where $h'$ is the constant-one prefix interpolation
for $S$, equivalently for $B_{n-1}(-rt)$. Since $0\le c<1$, the induction
hypothesis again gives the strict assertion outside the prefix.
Both new parameters have absolute value at most $r$. Antipodal
reflection commutes with $B_n(t)$, proving the suffix assertion.
\end{proof}

Work in signed coordinates $g=DG_\sigma$.
For a nonempty proper prefix $H$ let $U=X\setminus H$, and minimize
\[
             \tfrac12g^TQ_\sigma g
                 \quad\text{subject to }g_x\ge a\ (x\in H).      \tag{6}
\]
Let $q>0$ be the common row sum of $Q_\sigma$.
The equality-constrained candidate has $g_H=a\one$ and
\[
 g_U=a\one-aq(Q_\sigma)_{U,U}^{-1}\one,
 \qquad
 (Q_\sigma g)_H=aq\left[\one-
       (Q_\sigma)_{H,U}(Q_\sigma)_{U,U}^{-1}\one\right].          \tag{7}
\]
By (4), $Q_\sigma$ is a positive scalar times $B_d(\rho)$.
The suffix assertion of Lemma~\ref{lem:prefix} makes every entry of
the bracket strictly positive. Thus the candidate satisfies the
inequality KKT conditions with strictly positive multipliers.
Strict convexity proves that it is the unique constrained mode:
\emph{every observed constraint is active at every lexicographic stage}.

\begin{lemma}[Uniformly wrong lexicographic modes]\label{lem:mode}
At every lexicographic stage after the first, the mode satisfies
\[
       g^H_\rho(x_{|H|})
            \le -a\frac{\rho}{1+\rho+\rho^2}<0.                \tag{8}
\]
In particular it is at most $-a/4$ when $\rho\ge1/2$.
\end{lemma}
\begin{proof}
Put $L_n(x,y)=(-\rho)^{h(x,y)}$ and $C_n(t)=L_n(I+tP_n)$.
For odd $d$, the signed covariance, without its scalar factor,
is $C_d(-\rho)$. Split $C_n(t)$ by the first bit. Its diagonal block
is $C_{n-1}(-\rho t)$, its Schur complement is a positive scalar
times $C_{n-1}(\rho t)$, and the conditional mean of the second half
given constant one on the whole first half is
\[
                         c(t)\one,
              \qquad c(t)=\frac{t-\rho}{1-\rho t}.
\]
Consequently the prediction at the next point of a prefix obeys a
simple recursion: in the first half use the lower-dimensional
parameter $-\rho t$; in the second half use
$c(t)+(1-c(t))m'$, where $m'$ is the corresponding next-point prediction
with parameter $\rho t$. An empty lower-dimensional prefix has
prediction zero. This follows by the same block elimination as in
Lemma~\ref{lem:prefix}, now with $r=-\rho$; it is an algebraic identity,
not an application of that lemma with a negative kernel parameter.

Starting at $t=-\rho$, every encountered parameter is
$\pm\rho^p$, $p\ge1$, and a positive one has $p\ge2$.
Thus always $t\le\rho^2$, and
\[
 c(t)\le c(\rho^2)
       =-\frac{\rho}{1+\rho+\rho^2}<0,
\]
since $c$ is increasing. The recursion gives $m'\le0$;
any nonempty prefix encounters at least one second-half step.
It follows inductively that the next prediction is bounded by the
last display. Multiply by $a$, using the already proved activity
of all observed constraints, to obtain (8).
\end{proof}
The input permutation is therefore ordinary lexicographic order,
independent of both scale and width. The order has a simple description,
although it contains exponentially many distinct inputs.
For later use, (4) also implies a weaker assertion for any proper history:
stationarity at unseen coordinates and positivity of every precision
entry force at least one unseen mode coordinate to be negative.

\section{Confidence for the actual Gaussian posterior}
\begin{lemma}[Concentration around a constrained Gaussian mode]
Let a nondegenerate centered Gaussian on $\R^M$ with covariance
$\Sigma$ be conditioned on a full-dimensional convex polyhedron $C$.
Let $g_*$ minimize $V(g)=\frac12g^T\Sigma^{-1}g$ on its closure.
Then
\[
                 \E[\|g-g_*\|_2^2\mid C]
                                     \le M\lambda_{\max}(\Sigma). \tag{10}
\]
\end{lemma}
\begin{proof}
Integrate the divergence of $(g-g_*)e^{-V(g)}$ over the polyhedron,
first intersecting with a large ball. On each polyhedral face its
outward normal $n$ satisfies $(g-g_*)\cdot n\ge0$, because $g_*$
belongs to the closed convex region. The spherical truncation term
vanishes by Gaussian decay. Dividing by the conditioning probability gives
$\E[(g-g_*)\cdot\nabla V(g)\mid C]\le M$.
Optimality yields $(g-g_*)\cdot\nabla V(g_*)\ge0$ for feasible $g$.
Also
\[
 (g-g_*)\cdot[\nabla V(g)-\nabla V(g_*)]
       \ge\lambda_{\max}(\Sigma)^{-1}\|g-g_*\|_2^2.
\]
Combine the inequalities. Edges have zero surface measure; alternatively
use smooth convex approximation.
\end{proof}
Every diagonal entry of (3) is at most $\sigma^2$, so its largest
eigenvalue is at most $M\sigma^2$. Applying (10) to the signed posterior
and using Lemma~\ref{lem:mode}, for $\rho\ge1/2$, gives
\[
 \Pp\{g(x_t)\ge0\mid g(H_t)>a\}
                   \le16M^2\sigma^2/a^2\longrightarrow0.        \tag{11}
\]
This loose bound covers the full truncated Gaussian posterior.
Its mean differs from its mode by at most $M\sigma$ in Euclidean norm,
so its sign is wrong as well. A union bound over the $M-1$ stages proves
simultaneous fresh-sampling mistakes with probability tending to one.

\begin{lemma}[A convexly truncated Gaussian remains sub-Gaussian]
\label{lem:subgaussian}
If $G\sim N(0,\Sigma)$ is conditioned on a convex set $C$ of positive
probability, then for every vector $v$ and real $t$,
\[
 \E[e^{t(v^TG-\E[v^TG\mid C])}\mid C]
                         \le e^{t^2v^T\Sigma v/2}.            \tag{11a}
\]
Consequently if an unseen signed coordinate has posterior mean
$\mu<0$, its correct-label probability is at most
$\exp\{-\mu^2/(2\Sigma_{ii})\}$.
\end{lemma}
\begin{proof}
Completing the Gaussian square gives
\[
 \log\E[e^{tv^TG}\mid C]
   =\tfrac12t^2v^T\Sigma v+\ell(t)-\ell(0),\qquad
 \ell(t)=\log\Pp(G+t\Sigma v\in C).
\]
Prekopa's theorem \cite{prekopa} makes $\ell$ concave: it is the logarithm
of the integral of a jointly log-concave translated Gaussian times a
convex-set indicator. It is differentiable here, and differentiating
at zero gives $\ell'(0)=\E[v^TG\mid C]$.
The tangent bound $\ell(t)\le\ell(0)+t\ell'(0)$ proves (11a).
The final assertion is its optimized exponential Markov bound.
\end{proof}

\section{A fixed-scale consequence for sampling and posterior means}
The positive precision also gives a conclusion without a small-scale
limit. The input order in the next theorem may depend on the chosen
fixed scale and differs from the order in Theorem~\ref{thm:main}.

\begin{theorem}[Exponential errors at every fixed positive scale]\label{thm:fixed}
Fix odd $d\ge1$ and $\sigma>0$. There is a fixed input permutation
on which the Gaussian-process posterior mean predicts the wrong label
at every input after the first. At each of these rounds its posterior
sampling mistake probability is at least $1/e$.
For every sufficiently large finite neural width, the same permutation
has all $M-1$ posterior-mean errors and posterior sampling mistake
probability greater than $1/3$ at every designated round.
The latter conclusion holds at an explicit sufficient width
$N=\exp(O_\sigma(d))$; no such bound is asserted here for the mean sign.
\end{theorem}
\begin{proof}
Let $\mu=\E[g\mid g(H)>a]$. Integrating the Gaussian derivative in
an unconstrained coordinate gives $(Q_\sigma\mu)_i=0$ for $i\notin H$.
Every observed mean exceeds $a$ and every precision entry is positive.
Thus some unseen coordinate has $\mu_i<0$, by stationarity and the positivity of every coefficient. Select one, reveal its label, and repeat. This gives a deterministic
permutation for the fixed scale and makes every posterior-mean sign wrong.

The truncated Gaussian has a log-concave density. Marginalization and
integration over a half-line preserve log-concavity \cite{prekopa}.
Consequently an unseen coordinate $Y$ has a continuous distribution
function $F$ with $\log F$ concave. The probability integral transform
and Jensen's inequality give
\[
             \log F(\E Y)\ge\E\log F(Y)
                             =\int_0^1\log u\,du=-1.
\]
For the selected negative mean,
$\Pp(Y<0)\ge F(\E Y)\ge1/e$.
Conditional neural convergence below transfers this strict gap above
$1/3$, and transfers the strictly negative means. The final quantitative
subsection proves the displayed width order for the probability statement.
\end{proof}
In particular fresh sampling has expected mistakes greater than
$(M-1)/3$ at those finite widths. Since the sequence and true labels
are fixed, the draws can be independent; Hoeffding's inequality then
also gives an exponential-in-$M$ probability bound for deviations below
any smaller fixed fraction of mistakes. The theorem does not infer
a posterior-label majority error from a negative posterior mean.

\section{Transfer back to finite neural widths}
At fixed $\sigma>0$, the vectors
$a_j(\cos(b_j+w_j\cdot x))_{x\in X}$ are iid, with finite second
moments and nondegenerate covariance (3).
The central limit theorem gives $F_{N,\sigma}\Rightarrow G_\sigma$.
Every history region and its intersection with a query sign has Gaussian
boundary probability zero. Every history region has positive Gaussian
probability. Their probabilities therefore converge separately, and
division proves conditional correct-label probability convergence.
This accounts for posterior reweighting of incoming features;
unconditional covariance convergence alone would not establish it.

Choose $\sigma$ so that the sum of the Gaussian correct-label
probabilities is less than $\eta/2$. Then choose $N$ large enough that
the sum of the finitely many conditional-probability differences is
less than $\eta/2$. This works for every larger $N$.
Choosing the total bound also below $1/2$ gives majority errors at every
designated round. To transfer means as well, condition on the incoming
features before integrating the Gaussian readouts. Each output is a
centered Gaussian with variance at most $\sigma^2$, uniformly in $N$.
Thus the unconditional output tails are uniformly sub-Gaussian.
The fixed history probability converges to a positive number, so the
conditional outputs are uniformly integrable. Weak conditional convergence
therefore implies convergence of their means, preserving every fixed
strictly negative sign once $N$ is large enough.

\subsection{An explicit width bound for conditional probabilities}
This bound handles the potentially rare positive-margin conditioning
event. Write
\[
 A=\tfrac12\sigma^2(1-\rho)^{d+1},\qquad
 B=\tfrac12\sigma^2(1+\rho)^{d+1}.
\]
These are the smallest and largest eigenvalues of (3), by its tensor
spectrum and odd $d$. Every prefix event has probability at least
\[
 p_*:=\left(\frac A{2\pi B}\right)^{M/2}
               \exp\left[-\frac{M(a+\sqrt A)^2}{2A}\right].     \tag{12}
\]
Indeed its signed Gaussian density can be integrated on
$[a,a+\sqrt A]^M$, using determinant at most $B^M$ and inverse
covariance norm at most $A^{-1}$.

For $0<\delta\le10^{-3}$ set
\[
 R=3\left[\tfrac M2\log2+\log\frac1{\delta p_*}\right],
 \qquad \varepsilon=\frac{\delta}{4(M+R)}.
\]
Conditional on incoming features, the neural output is Gaussian with
empirical covariance $\Sigma_N$. Each covariance entry is an average
of independent variables in $[-\sigma^2,\sigma^2]$.
Hoeffding and an entrywise-to-operator bound imply
\[
 \Pp\{\|\Sigma_N-\Sigma_\sigma\|_{\rm op}>\varepsilon A\}
   \le2M^2\exp\left[-\frac{N\varepsilon^2 A^2}{2M^2\sigma^4}\right].
                                                               \tag{13}
\]
Thus the bad-covariance probability is at most $\delta p_*$ if
\[
 N\ge \frac{2M^2\sigma^4}{\varepsilon^2 A^2}
                                  \log\frac{2M^2}{\delta p_*}.  \tag{14}
\]

On the good event,
$(1-\varepsilon)\Sigma_\sigma\preceq\Sigma_N
\preceq(1+\varepsilon)\Sigma_\sigma$.
On the ellipsoid $V(x)=x^T\Sigma_\sigma^{-1}x/2\le R$, determinant
and quadratic-form comparison bound the density ratio between
$e^{-\delta}$ and $e^\delta$.
Both Gaussian laws have outside probability at most $\delta p_*$:
the elementary chi-square exponential moment gives the upper bound
$2^{M/2}e^{-R/3}$, uniformly for relative covariance error at most $1/2$.
Hence, for every measurable event $E$, its conditional-good-covariance
probability lies between
$e^{-\delta}(\Pp_{\rm GP}(E)-\delta p_*)$ and
$e^\delta\Pp_{\rm GP}(E)+\delta p_*$.
Mixing over the good and bad incoming-feature draws and dividing by a
history probability $p_H\ge p_*$ gives
\[
 \Pp_N(\text{correct}\mid H)
 \le\frac{e^\delta\Pp_{\rm GP}(\text{correct}\mid H)+2\delta}
              {(1-\delta)e^{-\delta}(1-\delta)}.        \tag{15}
\]
We used $\delta p_*\le\delta$ in the denominator.
For $\delta=10^{-3}$ and Gaussian correct probability at most
$1-1/e$, the right side is less than $2/3$.
At fixed $\sigma>0$, $A^{-1},B,M$, and $-\log p_*$ are all
$\exp(O_\sigma(d))$. The same is true of (14), proving the claimed
width order. The constants may be very large.

\begin{theorem}[A finite-parameter confident-mistake family]
There is an absolute constant $C$ such that, for every odd $d\ge1$
and $0<\eta<1/2$, the lexicographic construction works with
\[
 \sigma^2=\frac{a^2\eta}{64\,2^{3d}},\qquad
 N\ge\left\lceil\exp\{C(d+1)[d+\log(1/\eta)]\}\right\rceil.       \tag{15a}
\]
At these finite parameters, posterior-label majority is wrong at all
$2^d-1$ designated rounds, and fresh posterior sampling is wrong at
all of them with probability at least $1-\eta$.
No quantitative posterior-mean conclusion is added in this statement.
\end{theorem}
\begin{proof}
The chosen scale has $\rho>1/2$. By (11), every Gaussian correct-label
probability is at most $\eta/(4M)$. Use the conditional comparison
(15) with $\delta=\eta/(1024M)<10^{-3}$. Its right side is less than
$\eta/(2M)$, so a union bound and the majority assertion follow whenever
(14) holds. To bound that threshold, note that
$\sigma^2\le1/2$ and $\sigma^2\le1-e^{-2\sigma^2}\le2\sigma^2$.
Thus
\[
 \log(1/A)=O\bigl((d+1)[d+\log(1/\eta)]\bigr),\qquad B\le1.
\]
Formula (12) gives
$\log(1/p_*)\le C_0M[1+A^{-1}+\log(1/A)]$.
It follows that $\log R$ and $\log(1/\epsilon)$ have the same displayed
order. Taking logarithms in (14) proves (15a), after increasing an
absolute constant. This uses the explicit history-probability bound;
it does not substitute a dimension-dependent scale into a fixed-scale
big-O statement without controlling its constants.
\end{proof}

\begin{theorem}[Half the cube at a polynomially small prior scale]
There is an absolute constant $C$ such that for every odd $d\ge1$ and
$0<\eta<1/2$, at
\[
 \sigma^2=\frac{a^2}{64[d+1+\log(1/\eta)]},\qquad
 N\ge\left\lceil\exp\left\{C(d+1)\log\frac{d+1}{\eta}\right\}\right\rceil,
                                                               \tag{15b}
\]
there is a deterministic input order with one initial label on which
posterior-label majority is wrong at the final $2^{d-1}$ rounds, and
fresh posterior sampling is wrong at all of those rounds with probability
at least $1-\eta$. The order is the mean-selected order of Theorem~\ref{thm:fixed},
not asserted to be lexicographic. For fixed $\eta$ the sufficient width
is $\exp(O(d\log(d+1)))$ and the prior variance is of order $1/d$.
\end{theorem}
\begin{proof}
At every proper history select an unseen coordinate of smallest Gaussian
posterior mean. As in the proof of the fixed-scale theorem, stationarity
and positive precision imply
\[
 \min_{i\in U}\mu_i\le-a\frac{|H|}{|U|}
                         \frac{\min Q_\sigma}{\max Q_\sigma}.
\]
The selected scale has $\rho^{d+1}>1/2$. At each of the final $M/2$
rounds, $|H|\ge|U|$, so the chosen mean is at most $-a/2$.
Since $\Sigma_{ii}\le\sigma^2$, Lemma~\ref{lem:subgaussian} gives
\[
 p_{\rm GP}(\mathrm{correct}\mid H)
       \le e^{-a^2/(8\sigma^2)}
       =e^{-8[d+1+\log(1/\eta)]}\le\frac\eta{4M}.
\]
Use (15) with $\delta=\eta/(1024M)$. It makes each finite-network
correct-label probability less than $\eta/(2M)$ once (14) holds.
A union bound proves the sampling assertion, and each probability is
below one half. The deterministic histories contain the true labels,
so the chosen order does not depend on the learner's random draws.

For the width, the same estimates used after (15a) now give
$\log(1/A)=O((d+1)\log[d+1+\log(1/\eta)])$.
The extra $\log(1/\delta)=O(d+\log(1/\eta))$ terms in (14) are absorbed
by $O((d+1)\log((d+1)/\eta))$. This proves (15b) directly from the
finite-width comparison. No interchange of width and scale limits
is used.
\end{proof}


\begin{theorem}[Exponentially many confident errors at every fixed scale]
\label{thm:fixed-confident}
Fix any $\sigma>0$, odd $d\ge1$, and $0<\eta<1/2$. Put
\[
 S_{d,\sigma}=\frac{(1+\rho)^{d+1}}{1+\rho^{d+1}},\qquad
 T_{d,\sigma,\eta}=\left\lfloor
 \frac{a S_{d,\sigma}}{a+\sigma\sqrt{2\log(4M/\eta)}}
                         \right\rfloor,
 \quad \rho=e^{-2\sigma^2},\ M=2^d.                    \tag{15c}
\]
At every sufficiently large finite width, the mean-selected order
has wrong posterior-label majority at all its final $T_{d,\sigma,\eta}$
rounds, and fresh posterior sampling is wrong on all of them with
probability at least $1-\eta$. A sufficient width is
$N\ge\lceil\exp\{C_\sigma[d+\log(1/\eta)]\}\rceil$.
For each fixed $\sigma,\eta$,
\[
 T_{d,\sigma,\eta}
       =\Theta_{\sigma,\eta}\left(
               \frac{(1+e^{-2\sigma^2})^d}{\sqrt d}\right)
               \qquad(d\to\infty\text{ through odd values}).   \tag{15d}
\]
Thus exponential confident-majority failure does not require the prior
scale to shrink with dimension. The number in (15c) may be zero in
small dimensions; the assertion then has no designated tail rounds.
\end{theorem}
\begin{proof}
Use the exact row sum rather than the smallest precision entry.
Let $q$ be the common row sum of $Q_\sigma$ and $b$ its largest
entry. Formula (4) gives $q/b=S_{d,\sigma}$: the row sum of
$K_\rho+\rho P K_\rho$ is $(1+\rho)^{d+1}$, while its largest
entry is the diagonal value $1+\rho^{d+1}$.
For a history $H$, let $u=|U|>0$, $m=\min_{j\in U}\mu_j$,
and fix any $i\in U$. Stationarity and $\mu_j\ge a$ on $H$ give
\[
 0=\sum_jQ_{ij}\mu_j
   \ge a\sum_{j\in H}Q_{ij}+m\sum_{j\in U}Q_{ij}.
\]
Since $0<\sum_{j\in U}Q_{ij}\le ub$, division yields
\[
                   m\le-a\left(\frac{S_{d,\sigma}}u-1\right).
                                                               \tag{15e}
\]
Whenever $u\le T_{d,\sigma,\eta}$, the selected mean is at most
$-\sigma\sqrt{2\log(4M/\eta)}$. Lemma~\ref{lem:subgaussian}
therefore bounds its Gaussian correct-label probability by
$\eta/(4M)$. Apply (15) with $\delta=\eta/(1024M)$ and union bound
over the at most $M$ designated rounds. The finite-network correct
probabilities are below $\eta/(2M)$, proving both claims.

At fixed $\sigma$, $\log(1/A)=O_\sigma(d)$,
$-\log p_*\le\exp(O_\sigma(d))$, and
$\log(1/\delta)=O(d+\log(1/\eta))$. Taking logarithms in (14)
gives the stated sufficient width. For fixed $\sigma,\eta$, the
numerator of (15c) has order $(1+\rho)^d$, its denominator has order
$\sqrt d$, and its unfloored value tends to infinity. This proves (15d).
The same finite widths can satisfy the all-round sampling lower bound
of Theorem~\ref{thm:fixed}; only the final tail is claimed confident here.
\end{proof}

\section{A necessary exponential width and a narrow-width comparison}
\begin{theorem}[A finite-width evidence bound]
For each fixed $\sigma>0$, there is $C_\sigma<\infty$ such that,
at every neural width $N\ge1$ and on every input order labeled by
parity, the cumulative expected posterior-sampling mistakes are at most
\[
                  C_\sigma[d\log(d+1)+N\log(N+1)].      \tag{16}
\]
For $0<\sigma\le1$, an absolute-constant bound is
\[
 C\left[\frac{d+N}{\sigma^2}
             +d\log(d+1)+N\log(N+1)\right].                    \tag{16a}
\]
Posterior-label majority mistakes are bounded by twice these quantities,
and $2/3$ confidence failures by three times them.
Thus width one has polynomial lifetime mistake bounds.
Conversely the fixed-scale exponential sampling conclusion above
requires $N=\Omega_\sigma(2^d/d)$ as $d\to\infty$ through odd values.
\end{theorem}
\begin{proof}
Put the first neuron's incoming weights within $1/(16d)$ of $\pi/2$
and its bias within $1/16$ of the reduced teacher phase. Its signed
cosine is at least $\cos(1/8)>0.9$ everywhere.
Put its readout in $[2\sqrt N,3\sqrt N]$.
For all other neurons require only $|a_j|\le1/(4\sqrt N)$, leaving
their incoming parameters unrestricted. Since cosine is bounded by
one, the total normalized interference is less than $1/4$.
This parameter event has the correct margin greater than $1/2$ on
the entire cube. Its Gaussian probability is at least
\[
               \exp\{-C_\sigma[d\log(d+1)+N\log(N+1)]\}.
\]
The first incoming box costs $O_\sigma(d\log(d+1))$, its amplified
readout costs $O_\sigma(N)$, and each inactive readout interval costs
$O_\sigma(1+\log N)$. These estimates hold uniformly for $N\ge1$. For $\sigma\le1$, the incoming
boxes cost at most $C[d/\sigma^2+d\log(d+1)]$. The first readout costs
at most $CN/\sigma^2$, and each inactive readout costs at most
$C[1+\log_+(\sigma\sqrt N)]\le C\log(N+1)$.
Here $\log_+(u)=\max\{\log u,0\}$; use a subinterval of length
$\min\{\sigma,1/(4\sqrt N)\}$ in each inactive box.
This proves the explicit version (16a).

At each round the posterior wrong-sign probability is at most the
probability that the next margin observation fails. If $p_t$ is the
conditional probability of satisfying that margin, then
$q_t\le1-p_t\le-\log p_t$.
The sum telescopes to minus the log prior probability of the complete
observed history, which is at least the displayed global teacher-box
probability. This proves (16), including adaptive orders chosen before
the fresh prediction randomness. A majority error has $q_t\ge1/2$;
a $2/3$ confidence failure has $q_t\ge1/3$.
Finally an expected error count greater than $(2^d-1)/3$ forces
$d\log(d+1)+N\log(N+1)=\Omega_\sigma(2^d)$, hence
$N=\Omega_\sigma(2^d/d)$.
\end{proof}
In particular Theorem~\ref{thm:fixed-confident} gives an exponential
confident-majority separation from width one at every fixed positive
prior scale, as well as the all-round sampling separation.
Together with (14), this locates the width needed for the exponential
fixed-scale sampling failure between two exponential scales in dimension.
Equation (16a) also shows that, for each fixed $\eta$,
the polynomial-scale family (15b) needs
$N=\Omega_\eta(2^d/d)$ to have its stated exponential mistake count:
if $N<2^d$, its bound is $O_\eta(d^2+Nd)$, and if $N\ge2^d$ the
claim is immediate. In contrast, width one has $O_\eta(d^2)$ lifetime
expected mistakes at that same scale. The finite-width bound diverges
as $\sigma\downarrow0$ and does not reverse the limits in
Theorem~\ref{thm:main}.

\section*{Scope and verification}
The preceding independent cosine note \cite{sign-only} already gives
an exact product-kernel calculation for sign-only majority mistakes.
The refinements here are the Gaussian-bias precision identity,
positive observed margins, a fixed order across the common small-scale
limit, and confidence for the truncated Gaussian posterior.
The Gaussian-process limit itself is classical; see \cite{lee}.
No blanket novelty claim is made.
The full theorem is not Lean formalized or independently reviewed.
Finite exact checks in the accompanying payload verify matrix and KKT
identities rather than the analytic limit theorem.

\begin{thebibliography}{9}
\bibitem{prekopa} A. Pr\'ekopa,
\emph{On logarithmic concave measures and functions},
Acta Scientiarum Mathematicarum 34 (1973), 335--343.
\url{https://acta.bibl.u-szeged.hu/14411/1/math_034_335-343.pdf}.
\bibitem{sign-only} ChatGPT-6 Astra,
\emph{A one-neuron target that defeats wide sign-only Bayes},
independent research draft, 18 November 2026.
\url{https://github.com/kaarelh/AI-works/tree/main/neural-network-generalization/independent-research/chatgpt-6-astra-2026-11/papers/sign-only-fourier-counterexample}.
\bibitem{lee} J. Lee, Y. Bahri, R. Novak, S. S. Schoenholz,
J. Pennington, and J. Sohl-Dickstein,
\emph{Deep neural networks as Gaussian processes}, ICLR 2018.
\url{https://arxiv.org/abs/1711.00165}.
\end{thebibliography}
\end{document}
